\chapter{Termos de consentimento e autorização}
\label{ap:termos}

Este apêndice reúne os instrumentos utilizados para informar os participantes da pesquisa sobre seus direitos, obter o consentimento livre e esclarecido e formalizar autorizações específicas quando aplicáveis. Apresenta-se, em primeiro lugar, o modelo do Termo de Consentimento Livre e Esclarecido elaborado para esta pesquisa. Os modelos institucionais da UNISINOS (Termo de Confidencialidade, Termo de Autorização de Uso de Imagem e Autorização de Uso de Obra Fotográfica), aplicáveis conforme o caso, serão anexados nas versões finais a serem assinadas em campo e estão disponíveis no repositório do projeto.

\section*{Termo de Consentimento Livre e Esclarecido (TCLE)}

\noindent\textbf{Título da pesquisa:} Cardápios digitais com autoatendimento em bares e restaurantes de Porto Alegre: um estudo de caso sobre adoção, motivações e desafios pós-2020.

\noindent\textbf{Pesquisador responsável:} César Hoffmann, aluno do curso de Bacharelado em Ciência da Computação da Universidade do Vale do Rio dos Sinos (UNISINOS).

\noindent\textbf{Orientador:} Prof.\ Dr.\ Guilherme Lacerda --- Escola Politécnica, UNISINOS.

\noindent\textbf{Instituição:} Universidade do Vale do Rio dos Sinos (UNISINOS).

\noindent\textbf{Contato do pesquisador:} \emph{(a preencher: e-mail institucional e telefone para contato.)}

\subsection*{1. Objetivo da pesquisa}

Esta pesquisa tem por objetivo analisar como bares e restaurantes de Porto Alegre que adotaram cardápios digitais com autoatendimento desde 2020 percebem os benefícios, riscos e desafios técnicos e operacionais dessa transição, bem como a forma com que seus clientes vivenciam essa nova interação.

\subsection*{2. Procedimentos}

A participação consistirá em uma das seguintes atividades, conforme o perfil do participante:
\begin{alineas}
	\item \emph{Entrevista semiestruturada} com duração estimada de 45 a 60 minutos, conduzida presencialmente ou por videoconferência, gravada em áudio para posterior transcrição;
	\item \emph{Resposta a questionário on-line} de aproximadamente 3 a 5 minutos, hospedado no Google Forms, sem identificação nominal obrigatória.
\end{alineas}

\subsection*{3. Riscos, desconfortos e benefícios}

Os riscos associados à participação são mínimos e limitam-se à possibilidade de desconforto ao discutir aspectos operacionais ou comportamentais. Não há benefícios financeiros diretos. O benefício esperado é a contribuição para o avanço do conhecimento sobre transformação digital no setor de food service e para a melhoria das soluções tecnológicas disponíveis ao setor.

\subsection*{4. Garantias ao participante}

Ao participante são asseguradas as seguintes garantias:
\begin{alineas}
	\item participação voluntária, sem qualquer ônus ou prejuízo em caso de recusa;
	\item direito de retirar o consentimento a qualquer momento, sem necessidade de justificativa;
	\item sigilo quanto à identidade e a quaisquer informações que possam identificá-lo, com uso exclusivo de pseudônimos ou códigos nas publicações;
	\item acesso aos resultados consolidados da pesquisa mediante solicitação.
\end{alineas}

\subsection*{5. Tratamento de dados pessoais (LGPD)}

O tratamento de dados pessoais nesta pesquisa observa a Lei Geral de Proteção de Dados Pessoais (Lei n.\ 13.709/2018). Os dados coletados são utilizados exclusivamente para fins acadêmicos vinculados ao Trabalho de Conclusão de Curso e a eventuais publicações dele decorrentes. As gravações de áudio e os arquivos brutos serão armazenados em ambiente de acesso restrito ao pesquisador e ao orientador, e descartados em até 24 meses após a defesa do trabalho. O participante poderá, a qualquer momento, solicitar o acesso, a correção ou a exclusão de seus dados.

\subsection*{6. Declaração de consentimento}

\noindent Declaro que li e compreendi as informações acima, que tive a oportunidade de esclarecer dúvidas com o pesquisador, e que consinto livremente em participar desta pesquisa nos termos descritos.

\vspace{2em}

\noindent\rule{0.45\textwidth}{0.4pt}\hfill\rule{0.45\textwidth}{0.4pt}\\
Assinatura do participante\hfill Assinatura do pesquisador\\[1ex]
Local e data: \rule{0.5\textwidth}{0.4pt}

\section*{Termos institucionais complementares}

Os seguintes modelos institucionais da UNISINOS serão aplicados conforme a situação exigir, em versões assinadas anexadas ao trabalho final:
\begin{alineas}
	\item \emph{Termo de Confidencialidade} --- a ser assinado com gestores de estabelecimentos que compartilharem informações sensíveis, como dados financeiros, contratos com fornecedores ou métricas operacionais;
	\item \emph{Termo de Autorização de Uso de Imagem} --- a ser assinado por participantes que figurarem em fotografias incluídas no trabalho;
	\item \emph{Autorização de Uso de Obra Fotográfica} --- a ser obtida quando forem utilizadas fotografias produzidas por terceiros (por exemplo, imagens de divulgação dos estabelecimentos).
\end{alineas}
